\documentclass[11pt]{article}
\usepackage{fontspec}
\usepackage[french]{babel}
\usepackage{amsmath}
\usepackage{amssymb}
\usepackage{amsthm}
\usepackage{mathrsfs}
\usepackage{tikz-cd}
\usepackage{enumitem}

\newcommand{\printedpage}[1]{\marginpar{\footnotesize p.~#1}}
\newcommand{\msp}[1]{\marginpar{\footnotesize [#1]}}
\newcommand{\Fl}{\mathrm{Fl}}
\newcommand{\M}{\mathcal{M}}

\begin{document}

\printedpage{--} % no page number printed on this page
% head: (none)

\begin{center}
\vspace*{5cm}
{\LARGE\sffamily LES DÉRIVATEURS}

\vspace{1.5cm}
{\large\textbf{ALEXANDRE GROTHENDIECK}}

\vspace{2.5cm}
{\large\textbf{Chapitre XIII}}

\vspace{0.5cm}
{\large\textbf{Catégories de modèles (1)}}
\end{center}

\clearpage

\printedpage{--} % no page number printed on this page
% head: (none)

\vspace*{6cm}
\noindent Ce texte a été déchiffré et transcrit en \LaTeXe{} par M. Künzer. Il a été édité par M. Künzer, J. Malgoire, et G. Maltsiniotis. La transcription est aussi fidèle que possible au manuscrit. Pour les quelques corrections évidentes, ou rares commentaires des éditeurs, ainsi que pour la numérotation originale des pages du manuscrit, les caractères de machine à écrire \texttt{[typewriter]} entre crochets sont utilisés. Un point d'interrogation entre crochets signifie que l'on n'est pas sûr du mot qui précède. Les index, la liste des axiomes, le glossaire, et les références bibliographiques à la fin du chapitre ont été compilés par les éditeurs. Les numéros des pages dans les index sont ceux de l'original.

\noindent Cette édition est provisoire. Les remarques, commentaires et corrections sont bienvenus. Envoyer un message à :

\noindent\textbf{maltsin@math.jussieu.fr}

\hfill G. Maltsiniotis\hspace*{2cm}

\clearpage

\printedpage{1}
% head: A. Grothendieck : Dérivateurs, Ch. XIII, Édition des Universités de Montpellier II et Paris VII
% foot: Version 25 janvier 2004

\msp{page 1}

\section*{1\quad Couples de Quillen $(\Phi, \Psi)$}

Je veux d'abord donner une variante affaiblie de la notion de triple de Quillen (``model category'' de Quillen). Ce ne sera pas une notion tout-à-fait autoduale -- ainsi il y aura \textit{deux} variantes affaiblies, duales l'une de l'autre.

Mais avant encore, un sorite sur les LLP (left lifting property) et RLP (right lifting property) (\footnote[1]{donner définition en forme, avec renvoi à \textsc{Quillen}.}). Soit $\M$ une catégorie quelconque (le plus souvent, $\M$ sera grande). On considère des parties
\[
\Phi,\ \Psi\ \subset\ \Fl(\M)
\]
et on définit des applications
\[
\boxed{\begin{array}{ccc}
\Phi & \longmapsto & \Phi_* \\[1ex]
\Psi & \longmapsto & \Psi^*
\end{array}}
\qquad
\Fl(\M) \longrightarrow \Fl(\M)
\]
\texttt{[plutôt} $\mathfrak{P}(\Fl(\M) \longrightarrow \mathfrak{P}(\Fl(\M))$\texttt{,}\newline\texttt{où pour tout ensemble} $E$\texttt{,} $\mathfrak{P}(E)$ \texttt{désigne l'ensemble des parties de} $E$\,\texttt{]} par % [sic] \mathfrak{P}(\Fl(\M) -> \mathfrak{P}(\Fl(\M)), a closing parenthesis is missing in the first term
\[
\begin{aligned}
\Phi_* \ &= \ \{\text{ensemble des } p \in \Fl(\M) \text{ qui ont la RLP par rapport à } \Phi\} \\
\Psi^* \ &= \ \{\text{ensemble des } i \in \Fl(\M) \text{ qui ont la LLP par rapport à } \Psi\},
\end{aligned}
\]
et ensuite
\[
\boxed{\begin{array}{ccc}
\Phi & \longmapsto & \widetilde{\Phi} \\[1ex]
\Psi & \longmapsto & \overline{\Psi}
\end{array}}
\qquad
\Fl(\M) \longrightarrow \Fl(\M)
\]
\texttt{[plutôt} $\mathfrak{P}(\Fl(\M) \longrightarrow \mathfrak{P}(\Fl(\M))$\texttt{]} par % [sic] \mathfrak{P}(\Fl(\M) -> \mathfrak{P}(\Fl(\M)), a closing parenthesis is missing in the first term
\[
\begin{aligned}
\widetilde{\Phi} \ &= \ (\Phi_*)^* \\
\overline{\Psi} \ &= \ (\Psi^*)_* \ .
\end{aligned}
\]

\msp{page 2}

Je résume les propriétés les plus immédiates de ces opérations dans une

\noindent\textbf{Proposition 1.}

\begin{enumerate}[label={}]
\item[a)] $\Phi \longmapsto \Phi_*$, $\Psi \longmapsto \Psi^*$ \textit{sont des applications décroissantes de} $\Fl(\M)$ \texttt{[plutôt} $\mathfrak{P}(\Fl(\M))$\texttt{]} \textit{dans lui-même, pour la relation d'ordre} $\subset$\textit{, et} $\Phi \longmapsto \widetilde{\Phi}$\textit{,} $\Psi \longmapsto \overline{\Psi}$ \textit{des applications croissantes.}

\printedpage{2}
% head: A. Grothendieck : Dérivateurs, Ch. XIII, Édition des Universités de Montpellier II et Paris VII
% foot: Version 25 janvier 2004

\item[b)] \textit{On a}
\[
\Phi\ \subset\ \widetilde{\Phi},
\]
\textit{et pour qu'on ait $\Phi = \widetilde{\Phi}$, il faut et il suffit que $\Phi$ soit de la forme $\Psi^*$. On dit alors que $\Phi$ est} \textup{clos à gauche} (\footnote[2]{Donner définition en forme.})\textit{. L'ensemble des parties de} $\Fl(\M)$ \textit{closes à gauche est stable par intersections quelconques} (\footnote[3]{de façon précise, $\bigcap_i(\Psi_i^*) = (\bigcup_i \Psi_i)^*$}), \textit{et $\widetilde{\Phi}$ est la plus petite partie close de} $\Fl(\M)$ \textit{qui contienne $\Phi$ et la plus grande partie $\Phi'$ de} $\Fl(\M)$ \textit{telle que $\Phi'_* = \Phi_*$ (ou telle que $\Phi'_* \supset \Phi_*$).}

\textit{Énoncés symétriques pour $\Psi$ : on a}
\[
\Psi\ \subset\ \overline{\Psi},
\]
\textit{on a $\Psi = \overline{\Psi}$ si et seulement si $\Psi$} \texttt{[est]} \textit{de la forme $\Phi_*$, on dit alors que $\Psi$ est} \textup{clos à droite.} \textit{L'ensemble des parties de} $\Fl(\M)$ \textit{closes à droite est stable par intersections quelconques} (\footnote[4]{de façon précise, on a $\bigcap_i(\Phi_{i\,*}) = (\bigcup_i \Phi_i)_*$}), \textit{et $\overline{\Psi}$ est la plus petite partie close à droite contenant $\Psi$, et aussi la plus grande partie $\Psi'$ de} $\Fl(\M)$ \textit{telle que $\Psi'^* = \Psi^*$ (ou telle que $\Psi'^* \supset \Psi^*$).}

\item[] \hspace*{-2.5em}\msp{page 3}

\item[c)] \textit{Les applications}
\[
\Phi \ \longmapsto\ \Phi_*, \qquad\quad \Psi\ \longmapsto\ \Psi^*
\]
\textit{établissent des bijections inverses l'une de l'autre entre l'ensemble des parties $\Phi$ closes à gauche de} $\Fl(\M)$\textit{, et l'ensemble des parties $\Psi$ closes à droite. C'est un anti-isomorphisme d'ensembles ordonnés. On a alors}
\[
(\Phi \cap \Phi')_* \ = \quad
\begin{array}{l}
\textit{partie close à droite de } \Fl(\M) \ \textit{``engendrée par } \Phi_*,\ \Phi'_*\textit{'',} \\
\text{i.e. \texttt{[la]} \textit{plus petite partie close à droite contenant} } \Phi_*, \\
\Phi'_*, \text{ i.e. } \overline{\Phi_* \cup \Phi'_*}
\end{array}
\]
\texttt{[pour} $\Phi$\texttt{,} $\Phi'$ \texttt{parties closes à gauche]}\textit{, plus généralement}
\[
(\bigcap_i \Phi_i)_* \ = \ \overline{\bigcup_i (\Phi_{i\,*})}
\]
\textit{pour toute famille $(\Phi_i)_{i \in I}$ de parties closes à gauche de} $\Fl(\M)$\textit{. Formules duales} \texttt{[pour des parties closes à droite]}
\[
\left\{
\begin{aligned}
(\Psi \cap \Psi')^* \ &= \ \widetilde{\Psi^* \cup \Psi'^*} \\
(\bigcap_j \Psi_j)^* \ &= \ \widetilde{\bigcup_j (\Psi_j^*)}\ .
\end{aligned}
\right.
\]

\item[d)] \textit{Si} $\Phi, \Psi \subset \Fl(\M)$\textit{, alors on a}
\[
\Phi \subset \Psi^* \iff \Psi \subset \Phi_*, \qquad\quad \textit{on écrit alors } \Phi \longleftrightarrow \Psi,
\]
\textit{cela signifie que $\Phi$ a la LLP pour $\Psi$, ou que $\Psi$ a la RLP pour $\Phi$.}
\end{enumerate}

\printedpage{3}
% head: A. Grothendieck : Dérivateurs, Ch. XIII, Édition des Universités de Montpellier II et Paris VII
% foot: Version 25 janvier 2004

\msp{page 4}

\noindent\textbf{Proposition 2.} \textit{Soit $\Phi$ une partie close à gauche de} $\Fl(\M)$\textit{. Alors $\Phi$ satisfait les conditions} a) b) c) \textit{ci-dessous :}

\begin{enumerate}[label={}]
\item[a)] \textit{$\Phi$ contient les isomorphismes, et est stable par composition, et par cochangement de base (}i.e. \textit{si on a}
$\vcenter{\hbox{\begin{tikzcd}
A \arrow[r, "i"] \arrow[d] & B \arrow[d] \\
A' \arrow[r, "i'"] & B'
\end{tikzcd}}}$
\textit{carré cocartésien dans $\M$, avec $i \in \Phi$, on a $i' \in \Phi$).}

\item[b)] \textit{$\Phi$ stable par facteurs directs (en tant que partie de l'ensemble des objets de} $\underline{\Fl}(\M)$\textit{).}

\item[c)] \textit{$\Phi$ ``stable par limites inductives ordinales'', dans le sens suivant : Soit $I$ un ensemble bien ordonné, $(A_i)_{i \in I}$ un système inductif dans $\M$, tel que les $u_{ji} : A_i \longrightarrow A_j$ soient dans $\Phi$, et $A_{i_0} \xleftarrow{\ \sim\ } \varinjlim\limits_{i<i_0} A_i$ si $i_0 \in I$ est un élément limite. Supposons que $A = \varinjlim\limits_{I} A_i$ existe. Alors $\forall\ i_0 \in I$, la flèche canonique $A_{i_0} \longrightarrow A$ est dans $\Phi$.}
\end{enumerate}

Il y a un énoncé dual pour une partie $\Psi$ close à droite; elle contient les isomorphismes, est stable par composition et par changement de base, par facteurs directs, et enfin stable par limites projectives ordinales.

\msp{page 5}

On aura une réciproque partielle dans les catégories accessibles. C'est pour l'énoncé de cette réciproque que j'ai mis à part la condition \textit{b}) de stabilité par facteurs directs, au lieu de la bloquer avec \textit{a}).

\noindent\textbf{Définition 1.}

\begin{enumerate}[label={}]
\item[\textit{a})] (Pour mémoire.) RLP, LLP, la relation $\Phi \longleftrightarrow \Psi$ entre parties de $\Fl(\M)$ (relation non symétrique!).

\item[\textit{b})] (Pour mémoire.) Partie de $\Fl(\M)$ close à gauche, close à droite.

\item[\textit{c})] Soit $\Phi \longleftrightarrow \Psi$. On dit que le couple $(\Phi, \Psi)$ a la \textit{propriété de factorisation de Quillen}, si \textit{toute} flèche $f$ de $\M$ se factorise en
\[
f = pi, \qquad\quad i \in \Phi,\ p \in \Psi.
\]
On dit que $(\Phi, \Psi)$ est un \textit{couple de Quillen} si c'est un couple associé par \texttt{[les]} propriétés LLP et RLP ($\Phi \longleftrightarrow \Psi$), et \texttt{[si]} elles \texttt{[les parties} $\Phi$ \texttt{et} $\Psi$\texttt{]} sont closes (\textit{i.e.} $\Phi = \Psi^*$, $\Psi = \Phi_*$), et si elles ont la propriété de factorisation de Quillen.
\end{enumerate}

J'énonce dès maintenant le résultat technique essentiel, assurant l'existence de beaucoup de couples de Quillen :

\end{document}
